\chapter{Reference: tables}
\label{ch:refman-table}

\begin{aside}
A table is a column of rows of columns. The expressive power
is in the helpers: column widths, header styling, zebra
stripes, alignment.
\end{aside}

\section{Working example}

\begin{lstlisting}[language=C++]
#include <maya/maya.hpp>
#include <vector>

using namespace maya;
using namespace maya::dsl;

struct TableDemo {
    struct Row { std::string name; int age; std::string job; };
    struct Model {
        std::vector<Row> rows{
            {"ada",    36, "mathematician"},
            {"alan",   42, "logician"},
            {"grace",  85, "compiler"},
            {"linus",  54, "kernel"},
        };
    };
    struct Quit {};
    using Msg = std::variant<Quit>;

    static Model init() { return {}; }
    static auto update(Model m, Msg)
        -> std::pair<Model, Cmd<Msg>> {
        return {m, Cmd<Msg>::quit()};
    }

    static Element view(const Model& m) {
        auto cell = [](std::string s, int w) {
            s.resize(w, ' ');
            return text(s);
        };

        auto header = h(
            cell("name", 10) | Bold,
            cell("age",   5) | Bold,
            cell("job",  16) | Bold
        );

        std::vector<Element> body;
        bool zebra = false;
        for (auto& r : m.rows) {
            auto row = h(
                cell(r.name,         10),
                cell(std::to_string(r.age), 5),
                cell(r.job,          16)
            );
            if (zebra) row = row | Bg<30, 30, 36>;
            zebra = !zebra;
            body.push_back(row);
        }

        return v(
            header,
            v(body),
            blank_,
            t<"q to quit"> | Dim
        ) | pad<1> | border_<Round>;
    }

    static auto subscribe(const Model&) -> Sub<Msg> {
        return key_map<Msg>({ {'q', Quit{}} });
    }
};

int main() { run<TableDemo>({.title = "table"}); }
\end{lstlisting}

\section{Notes}

\begin{itemize}[leftmargin=*]
  \item Pad cells to a fixed width before passing to
    \code{text}; uniform widths make rows align.
  \item Use \code{Bg<...>} for zebra stripes; readers track
    rows by the stripe.
  \item For sortable tables, sort the underlying vector;
    don't try to sort in render.
\end{itemize}

\section{Recap}

\begin{recap}
\begin{itemize}[leftmargin=*]
  \item Build cells, pad to fixed width.
  \item Header row above, body below.
  \item Zebra stripes with alternating background.
\end{itemize}
\end{recap}

\section*{Workshop}

\begin{enumerate}[leftmargin=*]
  \item Add a sort key; press a column letter to sort.
  \item Highlight the row under the cursor.
  \item Add a footer with the row count.
\end{enumerate}
